\chapter{Appendix}
\label{cha:Appendix}

\begin{lstlisting}[caption={Question–Answer generation prompt}, label={lst:questionAnswerGenerationPrompt}]
You are generating a scientific 2-hop question-answer-evidence (Q-A) Tuple.
You are given 5 candidate scientific papers in the following format: 
[
  {
    "ArXiv": string,
    "text": string,
  }
]

Your task is to select the best pair of papers and generate ONE question that requires connected reasoning across exactly TWO papers.

PAPER SELECTION
First, select the best pair of papers.
The selected pair must satisfy:
-Paper A provides an intermediate entity, method, dataset, variable, or result
-Paper B uses, evaluates, extends, contrasts with, explains, or depends on that intermediate element
-The final answer requires combining both papers
-Do NOT select papers that are only loosely related or redundant.

REQUIREMENTS
A valid question MUST: 
-require exactly two reasoning steps (2-hop)
-use exactly two supporting papers
-require sequential reasoning:
    -Step 2 must depend on the result of Step 1
    -If Step 1 is removed, Step 2 should not be solvable
-NOT be answerable from either paper alone
-be self-contained and unambiguous
-require combining information across both papers
-only use information that is present in the papers, do not invent hypothetical connections or applications. Only use existing connections.
-question cannot contain explicit references to the papers or its content such as "in this paper", "the proposed methods" or similar

AVOID
Do NOT generate:
-multi-part questions (e.g., "What is X and what is Y?")
-vague or generic questions (e.g., "implications", "advantages")
-literature summary questions
-questions answerable from a single paper
-questions where papers are only topically related but not logically connected
-Do NOT use external scientific knowledge, commonsense assumptions, or unstated domain knowledge, only use the information that is present in the papers

OUTPUT FORMAT
{
    "usedPapers" : [
        {
            "arXiv": <paperId>,
            "role": <bridgeEvidence / bridgeAnswer depending on the papers role> 
        },
        {
            "arXiv": <paperId>,
            "role": <bridgeEvidence/ bridgeAnswer depending on the papers role>
        }
    ],
    "rejectedPapers": [<paperId1>, <paperId2>, <paperId3>],
    "reasoning": {
        "step1": <Explain your reasoning for step1 on paper A>,
        "step2": <use step1 reasoning to derive or support the final answer>,
        "connectionExplanation": <explain why the steps are connected and why step2 requires the result of step1>,
    }, 
    "questionDraft" : "<one clear, self-contained question requiring 2-hop reasoning, this serves as a draft for you final question. You can but don't have to explicitely reference the papers here to help formulate a better question.>",
    "question" : <decontextualize the questionDraft, so that it contains no explicit reference to the papers or external figures. Every reference to a paper should instead directly name the concept, phenomenon, method or similar that you are referring to, briefly describing it if needed. The question cannot contain explicit references to the papers or its content such as "in this paper", "the proposed methods" or similar.>",
    "answerWithPaperReferences" : <long-form paragraph integrating Paper A and Paper B with citations like [Paper A], [Paper B]>,
    "answerWithoutPaperReferences": <long-form paragraph that answers the question, without referencing the papers, but directly incorporatestheir information, so that it is standalone understandable>,
    "isNotSingleHop": <explain why the question you generated is not single hop>
}

Do not deviate from this schema. Do not add any preciding information like ```json. Only Answer with the valid json
Paper Texts:
\end{lstlisting}


\begin{lstlisting}[caption={Dataset Generation Judgement prompt}, label={lst:Dataset Generation Judgement prompt}]
You are evaluating a scientific question-answer (Q-A) example.
Your goal is to assess whether it is a valid 2-hop, evidence-grounded scientific multi-hop question.
A valid example must:
-require combining information from exactly TWO papers
-involve dependent reasoning (Step 2 must require Step 1)
-not be decomposable into independent sub-questions
-be fully answerable from the provided evidence
-be self-contained and unambiguous

IMPORTANT:
-For each criterion, assign one label:
-GOOD = fully satisfies the criterion
-BORDERLINE = partially satisfies the criterion
-BAD = does not satisfy the criterion
-Be strict: only assign GOOD if the criterion is clearly and fully satisfied.

INPUT:
Question: {question} 
Answer: {answer} 
bridgeEvidencePaperText: {paper1Text}
bridgeAnswerPaperText: {paper2Text}
Reasoning Steps: {reasoningSteps}

EVALUATION CRITERIA
A. Reasoning Structure
Multi-hop Validity: Does answering the question require combining both papers?
Assesses whether answering the question requires integrating information from both papers. Thus, this criterion is only satisfied, if neither paper provides sufficient information for deriving the complete answer.
-GOOD: Both papers are strictly required, the question cannot be answered using only one paper
-BORDERLINE: Both papers contribute but one may be sufficient
-BAD: Only one paper is sufficient (single-hop)

Dependency Strength: 
evaluates whether the steps in the reasoning chain are sequentially dependent, in order for the conclusion to be drawn. Thus, the second reasoning step cannot be solved without first applying or interpreting the first one.
Step 2 must require the result, entity, method, dataset, variable, or conclusion obtained in Step 1.
GOOD: Step 2 strictly requires Step 1
BORDERLINE: Partial dependence
BAD: Steps are independent (disconnected reasoning)

Non-Decomposability: Is the question NOT decomposable into independent sub-questions?
Verifies that questions cannot be broken down into independent single-hop sub-questions. Instead, reasoning steps must be connected in such a way, that separately solving them would not be sufficient for deriving an answer.
-GOOD: Cannot be split; requires joint reasoning
-BORDERLINE: Partially decomposable
-BAD: Clearly decomposable into independent sub-questions

B. Evidence Grounding
Evidence Distribution: Are both papers required and non-redundant?
Examines whether both papers contribute distinct and necessary evidence, ensuring that required evidence is spread across both papers and questions cannot be answered by using a single dominant source, with redundant information.
-GOOD: Each paper contributes distinct, necessary information
-BORDERLINE: Some overlap or redundancy
-BAD: One paper is sufficient; the other is redundant

Answerability: Is the answer fully supported by the provided evidence?
Assesses whether the answer is fully supported and obtainable by only relying on information present in the  papers, without requiring external knowledge or speculative inference, not explicitly supported by either paper. 
-GOOD: Fully supported by cited evidence
-BORDERLINE: Partially supported
-BAD: Not supported or contradicts the evidence

C. Dataset Quality: 
Decontextualization
Is the question self-contained and unambiguous? The question cannot contain explicit references to the papers or its content such as "in this paper", "the proposed methods", " this approach" or similar.
-GOOD: Fully self-contained; all entities clearly defined
-BORDERLINE: Minor ambiguity
-BAD: Not understandable without external context

OUTPUT FORMAT
{{
    "multiHopValidity": {{
        "judgement": "<GOOD / BORDERLINE / BAD>",
        "explanation": "<explain your judgement for multiHopValidity>"
    }},
    "dependencyStrength": {{
        "judgement": "<GOOD / BORDERLINE / BAD>",
        "explanation": "<explain your judgement for dependencyStrength>"
    }},
    "nonDecomposability": {{
        "judgement": "<GOOD / BORDERLINE / BAD>",
        "explanation": "<explain your judgement for nonDecomposability>"
    }},
    "evidenceDistribution": {{
        "judgement":"<GOOD / BORDERLINE / BAD>",
        "explanation": "<explain your judgement for evidenceDistribution>"

    }},
    "answerability": {{
        "judgement": "<GOOD / BORDERLINE / BAD>",
        "explanation": "<explain your judgement for answerability>"

    }},
    "decontextualization": {{
        "judgement": "<GOOD / BORDERLINE / BAD>",
        "explanation": "<explain your judgement for decontextualization>"

    }},
    "confidence": <0 to 1 rate how confident you are in the judgements>
}}

Do not deviate from this schema. Do not add any preciding information like ```json. Only Answer with the valid json
\end{lstlisting}


\begin{lstlisting}[caption={DIRE prompt}, label={lst:DIRE prompt}]
You are given 2 scientific paper texts as well as a logical connection that exists between that papers. Your task is to determine wether that connection is explicitely mentioned in one of the 2 papers, or if is drawn based on their content, but not explicitely mentioned. 

IMPORTANT:
-Use only the information from the given papers
-Do NOT rely on external knowledge
-Be strict: if any critical information is missing, the answer is NOT recoverable
-Partial or speculative answers count as NOT answerable

INPUT:
paper1: {paper1}
paper2: {paper2}
connection: {connection}

OUTPUT FORMAT
{{
    "explanation" : "<describe your reasoning>",
    "isPresent": <true/ false depending on if the connection is explicitely mentioned in at least one of the papers, false if not>
}}
Do not deviate from this schema. Do not add any preciding information like ```json. Only Answer with the valid json
\end{lstlisting}



\begin{lstlisting}[caption={LLM cited text span extraction prompt}, label={lst:LLM cited text span extraction prompt}]
You will be given the text of a scientific paper and a claim which was generated based on the content of the paper. 
You can assume that all information in the paper is factually correct. Your job is, that for the given claim you should extract the section of the paper which provides the 
best factual basis for justifying the statement. For that you should at least extract one sentence and up to 5 sentences which must be consecutively present in the paper. 
The extraction can only contain sentences that are directly quoted as they are written in the document. You are not allowed to make any adjustments whatsoever to them.
Here is a list of criteria that defines what a good
extraction looks like:
    -The claim is the logical consequence of the facts in the context
    -The claim and the context cover the same topic
    -each topic that is present in the claim is also covered in the context
    -the context does not contradict the claim
    -the claim includes only information provided in the context or correct generalizations or interpretations of that information
    -the context includes every piece of information needed to justify each individual piece of the claim
    -the context and the claim convey the same meaning, even though the claim might be worded differently
    -the context includes every semantic entity that is present in the claim
    -the information in the context is unambiguous and its interpretation can only lead to the logical conclusions that are drawn in the claim
Search for the context, that fulfills all those criteria the best.

Input: 
-text of a scientific paper
-claim generated by an LLM on the basis of this paper

Output:
-The extracted context

Answer with only the extracted sentences, exactly as they are consecutively! written in the document
You are not allowed to hallucinate or make any adjustments to the sentences. You are also now allowed to add additional comments or headllines.
Strictly reply with only those sentences. 
Try to select as few sentences as possible. The response must include every information needed to fullfill the criteria, but should not contain any unnecessary information,
that is not needed to justify the claim.


Scientific Paper Text : "{documentText}"
Claim: "{answer}"
\end{lstlisting}
